\chapter{Language assistance, causal discovery and uncertainty}
\label{ch:la03-discovery}

This chapter follows a graph suggestion through the operations that could
make it useful for a decision. The opening comparison separates linguistic
knowledge from observations and interventions. The subsequent studies develop
information bounds, numerical alternatives, search failures, expert evidence
and executable graph contracts. Temporal and latent-variable examples then
show what those operations cannot determine from the supplied information.
The final identification study asks whether the resulting answer still
addresses the original population and action. These are successive obligations,
not interchangeable measures of discovery quality.

\section{The scientific object behind a graph proposal}

A banking analyst may begin with a list of columns and a request to explain why
balances move after a change in deposit pricing. A language model can recognize
the vocabulary, propose a market-rate variable, distinguish a customer offer
from a realized rate, and suggest that customer characteristics affect both
assignment and retention. These suggestions are useful because a statistical
algorithm does not recover the meaning of a column from its numerical values.
The scientific question nevertheless remains open. The proposed relation may
describe a mechanism in another population, a historical association, a direct
effect, an indirect path, or simply a familiar phrase. The role given to the
suggestion determines the claim the resulting graph can support.

The inspected literature offers several roles rather than one method called
``LLM causal discovery.'' Pairwise language judgments, breadth-first graph
construction, retrieval-assisted discovery, response-based independence
queries, statistical prompting, soft graph priors, multi-agent debate and
numerical experiment design alter different parts of an analysis
\citep{v12_la_2305_00050,v12_la_2402_01207,v12_la_2402_15301,
v12_la_2406_07378,v12_la_2402_01454,v12_la_2405_13551,
v12_la_2407_15073,v12_la_2608_09696}.
They should be compared by what information they consume, what object they
change, and what uncertainty survives that change.

Let $D$ denote numerical observations, $T$ the retrieved or supplied technical
text, and $M$ the language-model mechanism, including its version, prompts and
decoding policy. A proposal is then a function or random output of
$(D,T,M)$. It is not an additional independent sample from the banking process.
If the model proposes a graph $G$, the relevant questions are whether the
variables have the intended meaning, whether the graph restrictions are
scientifically warranted, whether the data support the inference procedure,
and whether the target effect is identified under those restrictions. These
questions remain distinct even when one conversational interface presents all
of them together.

The following analysis uses inspected primary methods and evaluations. Reported
benchmark gains are the authors' results; no external discovery package or
language-model endpoint is executed here. The finite examples and probability
calculations below are project constructions. Their role is to test the logic
of proposed interfaces and interpretations, not to claim bank validation.

\section{What language-only discovery measures}

The broad proof of concept of K{\i}c{\i}man and colleagues motivates language
support for several causal tasks, including the elicitation of background
knowledge and reasoning about causal relationships
\citep{v12_la_2305_00050}. A model can answer a question such as whether an
increase in a borrowing rate affects demand without receiving a numerical
sample. Such an answer may encode useful inherited knowledge. It cannot reveal
which of several observationally compatible assignment mechanisms generated a
particular bank's records unless the prompt or the model's knowledge provides
additional applicable information.

Feng and colleagues investigate relation occurrences in accessible OLMo and
BLOOM pretraining corpora and report that relation frequency is associated with
prediction success \citep{v12_la_2407_19638}. Their controlled additions of
synthetic relations and the effects of reversed or negated statements help
separate recall from a claim of general discovery. The scope matters: their
pretraining investigation concerns accessible models and predominantly
commonsense relationships. It is evidence for a contamination concern and a
test design, not proof that every model operates only by memorization.

A sensible banking evaluation therefore crosses two kinds of novelty. One can
change the variable names while holding the mechanism fixed, or change the
mechanism while retaining familiar labels. The first tests reliance on
surface vocabulary. The second asks whether a familiar story prevents the
system from responding to contrary evidence. Neither test should declare a
new label intrinsically uninformative: operational definitions often carry
real information about intervention feasibility and time order. Removing all
meaning from a variable name changes the task as well as its contamination risk.

Context must also be handled carefully. In the inspected reliability study,
positive and negative contexts alter affirmative response rates, with contexts
generated by a model and assessed by human annotators
\citep{v12_la_2407_19638}. Some contextual changes legitimately change a causal
claim. A policy that affects demand in one regulatory regime may have no effect
in another. A robustness test should either hold the target world fixed or
declare the changed world and its new truth. Scoring every changed answer as
instability would conflate semantic responsiveness with error.

Jiralerspong and colleagues use breadth-first organization to reduce graph
construction calls \citep{v12_la_2402_01207}. This offers an engineering
alternative to asking about every pair separately. Call complexity is not total
resource cost. A call that lists many nodes may have a growing prompt and
response, require retries, or depend on a long context. A complete comparison
must record input and output tokens, wall time, cache policy, model cost,
failure frequency and the review time needed to adjudicate proposals. Rejecting
a cycle is useful structural validation, but does not make the remaining
acyclic graph true.

\section{Response tests are not sampling-law independence tests}

Cohrs and colleagues investigate language-model judgments in a
constraint-based discovery workflow \citep{v12_la_2406_07378}. This is a
different mechanism from applying a conditional-independence test to observed
bank records. For example, the statement
$X\mathrel{\perp\!\!\!\perp}Y\mid Z$ is a claim about a distribution. A
model's frequency of saying ``YES'' to that statement is a property of a
prompted response process. Even a correctly estimated response probability
does not establish the distributional statement.

Static inspection of the pinned \texttt{causal\_GPT} source shows a further
issue that can be examined without making any model calls. Its response
aggregation counts YES and NO within the same response collection. The
inspected test forms proportions $\widehat p_N=N_N/n$ and
$\widehat p_Y=N_Y/n$, pools them as
$\widehat p=(N_N+N_Y)/(2n)$, and uses the denominator
\begin{equation}
 \widehat s_{\mathrm{ind}}=
 \sqrt{\frac{2\widehat p(1-\widehat p)}{n}}.
 \label{eq:la03-independent-denominator}
\end{equation}
The formula resembles an independent two-proportion standard error. Its
interpretation needs the actual sampling relationship between the counts.

Consider the deliberately simple project response model in which independent
draws take values NO, YES or OTHER, with probabilities $p_N,p_Y,p_O$. Let
$I_N$ and $I_Y$ be indicators for the first two categories in one draw. The
categories are mutually exclusive, so
\[
 \operatorname{Cov}(I_N,I_Y)=-p_Np_Y.
\]
Consequently the exact variance of the difference of the sample proportions is
\begin{equation}
 \operatorname{Var}(\widehat p_N-\widehat p_Y)
 =\frac{p_N+p_Y-(p_N-p_Y)^2}{n}.
 \label{eq:la03-multinomial-variance}
\end{equation}
This elementary calculation is a project derivation under the stated
independent categorical law. It is not a claim that actual model calls are
independent or stationary.

Under the binary symmetric response null $p_N=p_Y=1/2$, the true variance is
$1/n$, whereas the square of
Equation~\eqref{eq:la03-independent-denominator} is $1/(2n)$. The denominator
is too small by a factor $\sqrt{2}$. For $n=100$ and a nominal one-sided
level of $0.01$, summing the exact binomial probabilities over the inspected
formula's rejection region gives approximately $0.044313$. The calculation
enumerates all counts with their rational binomial probabilities; it does not
simulate responses or execute the downloaded implementation.

Correcting this variance would address one response-model issue. It would not
convert language judgments into a bank conditional-independence test. Real
calls may share context, cached outputs, decoding artifacts or a changing
backend. The source also converts the response decision to a numerical zero
or one at the independence-test interface. Those values need a type that
exposes their meaning. Passing them through an API named ``pvalue'' cannot
confer sampling-law calibration.

In a funding example, repeated model agreement that customer risk is unrelated
to pricing conditional on branch does not validate the bank's assignment law.
The remedy is to separate a language-evidence object from a numerical test
object. The former can carry source anchors, answer frequencies and unresolved
scientific questions. The latter must identify its observations, null,
sampling assumptions and calibration. Neither object should silently replace
the other.

\section{Retrieval, constraints and soft priors}

Retrieval changes the information available to a model. It may supply an
operational definition or an expert restriction that the numerical data cannot
identify. It can also supply irrelevant or contradictory evidence. Zhang and
colleagues' retrieval-assisted method reports mixed results: in the inspected
ASIA setting, documents reduce skeleton F1 from 1 to 0.73, while other settings
benefit \citep{v12_la_2402_15301}. The result argues for keeping retrieval
conditions visible in a comparison, rather than labeling retrieval as a
uniform improvement.

The reported orientation accuracy in that work is evaluated on recovered
true-positive edges. Its denominator excludes missing or false edges. For a
reader deciding whether a graph is useful for intervention, skeleton error and
orientation error both matter. A system must report the metric and its
denominator before drawing a conclusion about the full graph. Downstream
effect behavior remains a separate outcome.

A retrieved statement may concern an ancestral path rather than direct
adjacency. Ban and colleagues elicit ancestral constraints and use soft
penalties alongside numerical scores \citep{v12_la_2306_16902}. In a chain
$A\to M\to Y$, $A$ is an ancestor of $Y$ even when no direct $A\to Y$
edge exists. If an adapter turns ``can affect'' into ``directly causes,'' it
changes the scientific input. Path, adjacency, temporal-order and exclusion
constraints therefore need different representations and conversion tests.

Hard constraints delete graphs. If a wrong constraint removes the true graph,
increasing the numerical sample does not restore it within the constrained
search. Darvariu and colleagues investigate soft language-derived priors that
retain support for possible graphs \citep{v12_la_2405_13551}. In schematic
form, one can write a search objective as
\begin{equation}
 \widehat G_\lambda\in\arg\max_G
 \{S(D,G)+\lambda L(T,M,G)\},
 \label{eq:la03-soft-score}
\end{equation}
where $S$ is a numerical score, $L$ a language-derived preference and
$\lambda$ a declared strength. This notation is a common project comparison
device; the actual score and statistical assumptions must be taken from the
chosen method.

Positive prior support prevents one particular failure: a graph is not
logically excluded solely because a model dislikes it. A large wrong
preference can still dominate weak numerical evidence. Subsequent pruning can
again remove the relevant graph. The inspected experiments include selected
combined gains and cases where standalone priors harm performance. A useful
evaluation fixes the prior rule, compares $\lambda=0$, varies evidence quality
and preserves every failed setting. Tuning $\lambda$ after viewing the final
test graphs invalidates a simple interpretation of the reported comparison.

Takayama and colleagues instead feed statistical discovery information into
language prompts, then convert response probabilities into forced, forbidden
or unknown constraints for particular algorithms
\citep{v12_la_2402_01454}. The conversion differs across PC, exact search and
DirectLiNGAM. Their inspected examples contain both gains and failures and
caution against treating the procedure as an opaque automatic solution.
A bootstrap frequency, a response frequency and a posterior model probability
are three different objects. A user interface should not place them under one
unqualified confidence label.

Statistical prompting also raises a data-reuse issue. If $D$ determines the
initial graph summary and that summary shapes the language restriction, the
restriction is a function of $D$. A later fit cannot casually treat it as
independent external evidence. This does not make every such procedure
invalid; it changes the procedure that must be justified or evaluated.
Tracking data hashes and the full adaptive rule makes the reuse visible.
Sample splitting may help a particular construction, but it neither proves the
restriction true nor automatically supplies valid downstream uncertainty.

\section{Debate, disagreement and scientific evidence}

The multi-agent discovery approach of Le and colleagues assigns different
roles to language agents and uses their interaction to construct a graph
\citep{v12_la_2407_15073}. This can organize critique and expose alternative
stories. Its inspected assumptions exclude hidden confounding and cycles, and
the reported canonical benchmark scale is limited. Those conditions should
remain attached to a banking interpretation.

Several agents need not provide several independent pieces of evidence. They
may use the same model family, encounter overlapping training material, or
retrieve the same article. A majority can therefore amplify a common error.
Disagreement is useful as a review signal, but agreement cannot be converted
to a causal confidence interval without a model of the information and
dependence involved.

For servicing, one agent may propose that a contact causes distress, another
that distress causes contact, and a critic may resolve the conflict by choosing
the more plausible story. A more defensible workflow keeps both candidates,
asks which timestamp and assignment evidence would distinguish them, and
identifies the effects under each. A later scientific owner may supply an
applicable restriction. The agent's role is to formulate the unresolved
question and retrieve the evidence, not to manufacture the answer.

Candidate remedies include comparing one-agent and multi-agent conditions,
using the same numerical evidence and total budget, and clustering evaluation
results by shared mechanism rather than counting paraphrased prompts as
independent cases. These are proposed tests. No multi-agent trial has been
executed here, and diversity of models alone does not verify assumptions.

\section{Graph ambiguity and the target effect}

A graph matters because it changes the interpretation of an intervention.
Suppose an observed distribution $P$ is compatible with a class
$\mathcal G$. For a target $Q$, let $\Theta_G(P,Q)$ contain the effects
allowed by the distribution and assumptions under $G$. The scientifically
relevant target set is
\begin{equation}
 \Theta_{\mathcal G}(P,Q)=\bigcup_{G\in\mathcal G}\Theta_G(P,Q).
 \label{eq:la03-effect-union}
\end{equation}
This is a definition, not a claim that the set has been computed or that a
learned class covers the truth. An unrestricted member can make the union
unrestricted. A point estimate from the most fluent graph narrative would
answer a different, more strongly conditioned question.

The project's two-variable example makes this concrete. Under model F,
$A\sim\mathrm{Bernoulli}(1/2)$ and $Y=A$ with probability $0.9$, with an
independent flip otherwise. Under model R the roles are reversed:
$Y\sim\mathrm{Bernoulli}(1/2)$ and $A=Y$ with probability $0.9$. Both
models yield
\[
 P(0,0)=P(1,1)=0.45,\qquad P(0,1)=P(1,0)=0.05.
\]
Replacing $A$'s equation by an intervention gives an average effect of $0.8$
under F and zero under R. In this deliberately finite class the target set is
$\{0,0.8\}$. Neither acyclicity nor observational fit selects the effect.
A sufficiently informative scientific time order or intervention could remove
an alternative, but the example intentionally assumes only a coarse shared
recording time.

The distinction between set identification and confidence in a learned set is
equally important. Let $G_\star$ be the true graph, $\widehat{\mathcal G}$ a
random retained class, and $C_G$ a random set for the target under $G$. Define
$C=\bigcup_{G\in\widehat{\mathcal G}}C_G$. If the actual joint procedure
satisfies
\[
 \Pr(G_\star\notin\widehat{\mathcal G})\le\alpha_G,
 \qquad
 \Pr(\theta_\star\notin C_{G_\star})\le\alpha_\theta,
\]
then the elementary union bound gives
\begin{equation}
 \Pr(\theta_\star\notin C)\le\alpha_G+\alpha_\theta.
 \label{eq:la03-coverage-decomposition}
\end{equation}
The second premise must refer to the actual procedure; it cannot simply be
borrowed from a fixed-model interval after unaccounted model selection. No
independence between the two events is needed for the bound. Neither premise
has been established for arbitrary language-assisted temporal discovery here.
The formula exposes what would have to be supplied, rather than supplying it.

CPDAG and PAG representations encode ambiguity that should survive
serialization and explanation. Circle marks, bidirected relations and
unoriented edges cannot be replaced with arrows merely to produce a cleaner
diagram. A PAG is not an assertion that every latent variable has been named.
Similarly a finite list of proposed DAGs is not automatically a complete
equivalence class. The declaration validator preserves the marks and the
declared alternatives while delegating full class recognition and
identification to a formal engine.

\section{Temporal discovery is not longitudinal policy inference}

The word ``sequential'' refers to several distinct problems. Verma and
colleagues process successive batches of contemporaneous observations while
querying a language model as a noisy expert and retaining PAG information
\citep{v12_la_2506_16234}. The query budget and answer histograms concern
the information supplied to a discovery procedure. They do not identify the
value of a sequence of treatment decisions. Semantic entropy of model answers
does not by itself establish simultaneous graph coverage.

RealTCD combines language-derived initialization with numerical optimization
over temporal graphs and intervention targets
\citep{v12_la_2404_14786}. Its setup assumes causal sufficiency, a known
observational regime and an unrolled finite-lag DAG. The inspected theorem
conditions its population optimization result on model capacity, intervention
faithfulness, positivity and finite entropy. A stochastic optimizer with a
thresholded adjacency matrix requires additional analysis before it can inherit
any finite-sample guarantee.

The real cooling-system study does not have a complete ground-truth DAG. It
uses selected expert physical-location rules to assess edges. That is useful
domain evidence but a different evaluation from full known-truth graph
recovery. The method also models abnormal windows as intervention regimes.
In a bank, a sudden spread movement could reflect a common market shock, a
changed selection mechanism or measurement failure. Treating it as an
exogenous intervention requires substantive evidence, not a threshold crossing.

Kougioulis and colleagues' Large Causal Model is an adjacent numerical
transformer rather than a language-model wrapper
\citep{v12_la_2602_18662}. Its inspected scope includes additive-noise,
stationary, causally sufficient, lagged-only systems. These conditions can be
useful for a declared class; they exclude several concerns in private banking
states and common shocks. The external pretrained weights have not been
audited or executed locally. A language interface should not rename this
restricted temporal predictor as a general latent-graph solution.

For longitudinal servicing, the analyst instead needs a sequence of decisions,
available histories, treatment-confounder feedback, censoring and a value
target. A graph of temporal dependence is only one ingredient. The
intermediate account state can be affected by an earlier action while
confounding a later action's effect. The correct question and inferential law
must be specified before an agent chooses an estimator. None of the inspected
temporal discovery results closes the project's general latent temporal
effect-set problem.

\section{New interventions, mechanism proposals and information value}

Lampinen and colleagues show how passive learning from expert demonstrations
can support active causal strategies \citep{v12_la_2305_16183}. Their
constructive argument first acquires causal information through experiments,
then exploits it. The evaluated learner can intervene in its test environment.
This is a valuable qualification: passive training is compatible with active
information acquisition, but the new interventions supply information that an
analyst restricted to historical bank records may lack.

The inspected empirical settings use small finite acyclic systems and
constrained observations and actions. A banking translation must specify which
interventions are feasible, who can authorize them, their costs and their
effect on future information. Merely asking a model another question is not
an intervention on the customer or the funding market. The two actions have
different likelihoods, costs and ethical or operational constraints.

Murphy's Model Discovery Agent proposes executable mechanisms with an LLM and
fits and compares them numerically \citep{v12_la_2608_09696}. The inspected
studies include algebraic mechanisms, latent ODEs and stochastic latent
dynamics. The LLM expands the candidate library; BIC or SMC approximations
provide model evidence, and a model-disagreement surrogate approximates
expected information gain for experiment selection.

This architecture gives each component a clearer role than requesting a final
effect directly from the model. It also exposes new obligations. An adaptive
library may omit the true mechanism. Pruning can discard a mechanism that
would become useful later. BIC weights within a retained pool are conditional
model weights, not uniform graph-confidence statements. Residual checks can
reveal some misspecification but cannot certify that no policy-relevant
mechanism is absent.

The paper explicitly states that its finite held-out prediction tests do not
establish causal identification. The experiments use open-loop policies, and
the reported comparisons focus on sample efficiency rather than matched total
compute. The chemistry comparison retains matched evaluable cells from a
larger attempted collection. Those denominators must accompany the gains.
For a bank, the value of information is the expected improvement in an actual
decision after paying acquisition and implementation costs. Information gain
within a model pool can be a design surrogate, but it is not that monetary
quantity without an explicit connection to the bank's loss.

\section{An auditable discovery interface}

Zheng and colleagues describe causal-learn+ with agents assisting data
inspection, assumption elicitation, tool choice and interpretation around a
formal discovery core \citep{v12_la_2606_23608}. They emphasize separate
objects for suggestions, user decisions, algorithm settings and graph outputs.
The design is a useful default for a banking library because it makes changes
to scientific input visible. It is not a universal impossibility result for
hybrid discovery; a hybrid can be evaluated as its own fully specified method.

CausalAgent supplies another conversational architecture, including MCP tool
interfaces \citep{v12_la_2602_11527}. Static source review finds that an
inspected OLC tool calls a function not imported or defined in that module,
and a parser can return an empty dictionary after malformed JSON. These are
integration findings about pinned source, not measurements of a running
service. They show why transport success, parse success and scientific success
need separate dispositions. A successful model message is not a valid graph.

For each request, an auditable interface should retain the business target,
variable definitions and clocks, numerical-data identity, retrieved-source
identity, model and prompt version, proposal, applied restrictions, selected
algorithm and returned graph class. A reviewer should be able to distinguish
an expert's asserted exclusion from an algorithmic orientation and from a
model's suggestion. Changed restrictions create a new analysis identity;
they should not overwrite the prior result invisibly.

Failure recovery must preserve scientific meaning as well as computation. If a
job stops after producing a partial graph, replay must not publish it as a
complete answer. If a source or data hash changes, an earlier identification
claim cannot simply be reused. A report should expose a refused or unknown
effect when that is the formal outcome. The completed local operation contract
provides a bounded implementation surface; new discovery adapters still need
their own source, semantic and execution qualification.

\section{What would establish useful progress}

A meaningful comparison begins with a fixed question and information set.
The numerical baseline, language-only baseline, retrieval condition, hybrid
and agent-assisted workflow should receive explicitly matched information or
have their information advantage declared. Accuracy on familiar graphs is
only one outcome. The study also needs identification correctness, retention
of compatible alternatives, valid abstention, effect-set coverage where a
guarantee is claimed, interval or set width, decision loss and full cost.

Every request belongs in the denominator. Failed parsing, unavailable models,
timeouts, nonidentification and legitimate refusals have different meanings
and should be recorded separately. A graph metric conditional on successful
requests can be informative, but it must accompany the success probability
and the selection rule. A narrow accurate set emitted only in easy cases is
not the same product as one that delivers useful valid output reliably.

The proposed remedies are correspondingly bounded. Novel labels test one
shortcut; counterevidence tests retrieval dependence; correctly typed
constraints prevent a specific semantic conversion; fixed prompts and budgets
make an experiment reproducible; and preserved graph classes prevent one form
of overstatement. None of these repairs creates missing bank assignment
truth, resolves observational equivalence by itself, or proves arbitrary
dependent-sample calibration.

This opening comparison supports language assistance for explicit proposals
and reviewable workflow support, with each hybrid method qualified separately.
The following studies develop that position through information limits,
numerical alternatives, source-specific failures and decision consequences.
They distinguish the restricted results already available from the general
temporal graph and useful-precision questions that remain open. New scientific
information, justified assumptions or additional experiments may change the
answer; more fluent explanations of the same information cannot.
